\documentclass[11pt]{ctexart}
\input{D:/aaa_D_usual/DSH/111学习助手/.dsh/skills/lecture-deepdive/assets/theme-violet}
\usepackage{tikz}
\usetikzlibrary{arrows.meta}
\settags{第五章 · 杨氏双缝干涉精讲}{光学（大二）}

\begin{document}
\typeout{DIAG hsize=\the\hsize textwidth=\the\textwidth linewidth=\the\linewidth}

\doctitle{杨氏双缝：从光程差到云母片}%
{第五章 \S5.5--5.7 \ 双缝干涉与分波前装置 · 精讲讲义}%
{依据：\texttt{第五章 光的干涉.pdf}（83 页扫描版，逐页读图）p13--p28\ \textbar\ 2026-10-27}

\callout{\textbf{这份讲义解决什么。}你手上那张手写稿（云母片盖住一条缝、算出 $7.6\,\mu\mathrm{m}$）对应的不是习题册，而是课件本身：\textbf{p20 的 Ex.\,5.2，页脚印着 HOME WORK TWO}。整份作业只有这一问。
本讲义把 p13--p28 这 16 页串成一条能独立读的线：\S1 讲清"普通光凭什么也能干涉"；\S2 把光程与相位差接上；\S3 补完双缝的几何近似（课件在这里跳了一步，是全班最容易写错的地方）；\S4、\S5 逐题全解（\S5 就是你的作业，给了三条互相独立的路线）；\S6--\S8 讲对比度、空间相干性和三种分波前装置；\S11 把课件每一处跳步连页码带原话列出来；\S12 给三题自检。
讲义里凡标 \textbf{[补]} 的地方，都是课件没有写、由本讲义补出的完整推导；凡标 \textbf{[课件原话]} 的地方，都是直接引用那 16 页的句子。}

\flowerrule

% ---------------------------------------------------------------- §0
\hdr{\S0}{课件坐标、范围与记号}

\subhdr{0.1 这 16 页分别在讲什么}

本节取自 \texttt{第五章 光的干涉.pdf} p13--p28，共 16 页。p12 之前是相干条件的一般讨论，p29 起进入分振幅（薄膜）干涉，所以 p13--p28 刚好是一段自足的「分波前干涉」。

\tabtop
\noindent\begin{tabular}{@{}p{1.3cm}p{6.1cm}p{4.5cm}p{4.0cm}@{}}
\thead{课件页} & \thead{标题（原文）} & \thead{核心内容} & \thead{本讲义位置}\\
\tabmid
p13 & Interference of Two Ideal Point Sources & 两理想点源的场与强度 & \S1.2\\
p14 & Intensity Distribution of Two Ideal Point Sources & $I=I_1+I_2+2\sqrt{I_1I_2}\cos\delta$ 的极大极小 & \S1.2、\S2\\
p15 & Relationship Between Phase Difference and Optical Path Difference & $\delta=\frac{2\pi}{\lambda}\Delta L$，双曲线族 & \S2\\
p16 & Young's Interference --- Interference with Ordinary Light Sources & 单缝造点源、双缝切波前 & \S1.4、\S3\\
p17 & The Phase Difference in Young's Interference & 初相位相消、$R_2^2-R_1^2=2yd$ & \S3.1\\
p18 & Fringe Spacing in Young's Interference & $\Delta y=\lambda D/d$ & \S3.2\\
p19 & Example 5.1 Measuring the Wavelength of Light & 例题（有解） & \S4\\
p20 & Ex.\,5.2 Mica Sheet Covering a Slit & \textbf{作业 HOME WORK TWO}（无解） & \S5\\
p21 & The Physical Essence of Young's Interference & 空间周期性、把装置看成"放大器" & \S3.4\\
p22 & Contrast of Interference Fringes & $\gamma=(I_{\max}-I_{\min})/(I_{\max}+I_{\min})$ & \S6\\
p23 & Interference Intensity in Terms of Contrast & $\gamma=\frac{2A_1A_2}{A_1^2+A_2^2}$ & \S6\\
p24 & Intensity Distribution Expressed by Visibility & $I=I_0(1+\gamma\cos\delta)$ & \S6\\
p25 & Wavefront Division: Core Idea and Key Challenge & 波前分割的做法与"致命问题" & \S7\\
p26 & Fresnel Biprism（正文实为双平面镜） & $d\approx 2\alpha B$ & \S8\\
p27 & Fresnel Biprism & $d\approx 2(n-1)\alpha B$ & \S8\\
p28 & Lloyd's Mirror & $\Delta x=D\lambda/(2a)$、半波损失 & \S8\\
\tabbot

\subhdr{0.2 记号表（全文只此一套）}

\tabtop
\noindent\begin{tabular}{@{}p{2.2cm}p{13.0cm}@{}}
\thead{符号} & \thead{含义与单位}\\
\tabmid
$S$ & 单缝（点光源）所在位置；它到两缝的距离记为 $r_1$、$r_2$\\
$S_1,\,S_2$ & 两条缝；到屏上观察点 $P$ 的几何距离记为 $R_1$、$R_2$\\
$P(x,y)$ & 屏上观察点，$y$ 为离轴高度（向上为正），$x$ 为屏内横向坐标\\
$d$ & \textbf{双缝间距}（p16--p18 的用法）；单位 m\\
$D$ & 双缝到屏的距离；单位 m\\
$h$ & \textbf{云母片厚度}（p20 图上标的是 $d$，与缝距撞名；本讲义改记 $h$）；单位 m\\
$n$ & 云母片或棱镜的折射率；无量纲\\
$\lambda$ & \emph{真空}中的波长；单位 m（介质中波长见 \S2）\\
$k$ & 波数，$k=2\pi/\lambda$（p13、p17 的用法）；单位 $\mathrm{m^{-1}}$\\
$\delta$ & 两束光的相位差；单位 rad\\
$\Delta L$ & 光程差（optical path difference）；单位 m\\
$m$ & 级次，整数 $0,\pm1,\pm2,\dots$；\textbf{不是} $k$\\
$\gamma$ & 对比度（可见度），$0\le\gamma\le1$\\
$B,\,C$ & 双镜/双棱镜中：源到镜交点的距离 $B$，交点到屏的距离 $C$\\
$\alpha$ & 双镜两镜面夹角 / 双棱镜顶角；单位 rad\\
$a$ & 洛埃镜中源到镜面的高度\\
\tabbot

\warnbox{\textbf{记号撞车警告。}课件自己在两处用了同一个字母：p16--p18 的 $d$ 是\emph{缝距}，p20 图上的 $d$ 是\emph{云母片厚度}。这不是小事——你手写稿最后一行 $d=8\lambda/(n-1)$ 里的 $d$，和前面 $\Delta L\approx yd/D$ 里的 $d$ 已经不是同一个东西了。本讲义一律用 $h$ 表示片厚，就是为了让这两件事在纸上分得开。}

\flowerrule
% ---------------------------------------------------------------- §1
\hdr{\S1}{地基：叠加、强度、相干三条件}

读这一节要抓住一件事：\emph{相位差}是唯一的主角，其余全是它的配角。

\subhdr{1.1 叠加原理：先把"能相加"说清楚}

p3--p5 给的是场矢量相加：
\[
\vec{E}(P,t)=\vec{E}_1(P,t)+\vec{E}_2(P,t).
\]
\aside{成立条件是介质线性（p4 原话："When the principle of superposition holds in a medium, the medium is called a \textbf{linear medium}"）。强激光、非线性晶体里这条不成立，那时光会自己"干涉自己"产生新频率——本课不涉及。}

\subhdr{1.2 强度是时间平均：课件 p6 跳掉的一步 \textbf{[补]}}

p6 直接写出 $I(P)=I_1+I_2+2\sqrt{I_1I_2}\cos\delta(P)$。这一步不是"代入就完"，中间隔着一次时间平均，而时间平均正是"相干"与"不相干"的分水岭。完整走一遍。

单色波写成 $E=A\cos(\omega t-\varphi)$，探测器（眼睛、底片、CCD）响应的是长时间平均：
\[
I\equiv\langle E^2\rangle=\frac{A^2}{2}\langle 1+\cos 2(\omega t-\varphi)\rangle=\frac{A^2}{2}.
\]
\aside{$\langle\cdot\rangle$ 指对远大于 $1/\omega$ 的时间取平均；$\cos 2(\omega t-\varphi)$ 在一个周期内正负抵消，故只剩 $1/2$。这个 $1/2$ 后面还会出现，届时要能指回这里。}

两束光叠加，平方后有三项：
\[
E_1E_2=\tfrac{A_1A_2}{2}\Big[\cos\big((\omega_1-\omega_2)t-(\varphi_1-\varphi_2)\big)+\cos\big((\omega_1+\omega_2)t-(\varphi_1+\varphi_2)\big)\Big].
\]
第二项永远平均为零（快振荡）。第一项分两种命运：
\begin{itemize}
\item 若 $\omega_1\ne\omega_2$，它仍在振荡，平均为零 $\Rightarrow$ 交叉项消失，$I=I_1+I_2$；
\item 若 $\omega_1=\omega_2$，它变成常数 $\cos\delta$ $\Rightarrow$ 交叉项留下，得到 p6 的结果
\[
I=I_1+I_2+2\sqrt{I_1I_2}\,\cos\delta,\qquad \delta=\varphi_2-\varphi_1 .
\]
\end{itemize}

同一页还给出了等强度的特例（p14 原话 "If $I_1\approx I_2=I_0$"）：
\[
I=2I_0(1+\cos\delta)=4I_0\cos^2\frac{\delta}{2}.
\]
\aside{这里的 $I_0$ 是\emph{单束}光的强度，不是总强度；两束各为 $I_0$ 时总强度在 $0$ 与 $4I_0$ 之间摆动。}

\subhdr{1.3 相干三条件：同一条式子的三个开关}

p9--p12 把条件列成三条。它们不是三条独立的经验规则，而是上面那个交叉项的三个"归零开关"：

\tabtop
\noindent\begin{tabular}{@{}p{3.4cm}p{5.6cm}p{7.4cm}@{}}
\thead{条件（p12）} & \thead{不满足时哪一项归零} & \thead{来源}\\
\tabmid
频率相同 & $\cos\big((\omega_1-\omega_2)t-\Delta\varphi\big)$ 的平均为零 & \S1.2 展开式的第一项\\
振动方向平行 & 矢量点积 $\vec{E}_1\cdot\vec{E}_2=0$，$I=I_1+I_2$ & p11：垂直分量互不干涉\\
相位差恒定 & $\langle\cos\delta(t)\rangle=0$，交叉项平均掉 & p9："The phase difference must remain constant"\\
\tabbot

逐行读法：第一行的"平均为零"是时间平均（\S1.2）；第二行是\emph{空间}上的点积——两束光的振动方向差 $90^\circ$ 时，交叉项里的 $A_1A_2\cos 90^\circ=0$，不需要任何平均就直接没了；第三行是 $\delta$ 本身在漂，平均也归零。三条对应三种不同的"归零机制"，这就是为什么它们必须同时满足。

p8 还解释了第三行为什么对普通光源成立不了：原子自发辐射是一段段有限长的波列，初始相位随机跳变，所以两只独立的灯之间 $\delta(t)$ 随机游走。

\subhdr{1.4 杨氏的办法：让初相位自己约掉（p16、p17）}

p17 的两行式子是全章最漂亮的地方。设单缝处初相位为 $\varphi_0(t)$（它随时间乱跳），则
\[
\varphi_{10}=k(R_1+r_1)+\varphi_0(t),\qquad \varphi_{20}=k(R_2+r_2)+\varphi_0(t),
\]
相减：
\[
\delta=\varphi_{20}-\varphi_{10}=k\big[(R_2-R_1)+(r_2-r_1)\big].
\]
$\varphi_0(t)$ 出现了两次，直接消掉——这正是 p17 的原话："It is independent of the initial phase $\varphi_0(t)$."
再把单缝放在对称轴上，$r_1=r_2$（p17 图上的构型），于是
\[
\delta=k(R_2-R_1).
\]
\aside{p16 的装置是"单缝 $S$ $\to$ 双缝 $S_1,S_2$ $\to$ 屏"。单缝的作用不是"提供能量"，而是把一只扩展光源切成一个足够小的\emph{相干点源}；双缝再把\emph{同一波前}切成两块。这就是"分波前"（wavefront division）四个字的意思，p25 会给正式定义。}

最小数字例子（只为看清机制，不是作业）：$\lambda=500$ nm、$d=1$ mm、$D=1$ m、$y=1$ mm，则
\[
\Delta L\approx\frac{yd}{D}=\frac{10^{-3}\times10^{-3}}{1}=10^{-6}\ \mathrm{m}=2\lambda ,
\]
在该点两束光的光程差正好是 2 个波长，是明纹。

\flowerrule
% ---------------------------------------------------------------- §2
\hdr{\S2}{光程与相位差：为什么介质里要乘 $n$}

\subhdr{2.1 光程的定义（p15 的前提，p15 没证）\textbf{[补]}}

光在折射率 $n$ 的介质中传播，频率不变、波长变短：
\[
\lambda_n=\frac{\lambda_0}{n},\qquad k_n=\frac{2\pi}{\lambda_n}=\frac{2\pi n}{\lambda_0}=nk_0 .
\]
走几何长度 $s$，相位滞后为
\[
k_n s=\frac{2\pi}{\lambda_0}\,(ns)\equiv\frac{2\pi}{\lambda_0}L,\qquad L\equiv ns .
\]
$L=ns$ 就叫\textbf{光程}：把"在介质里走了 $s$"折算成"在真空中要走 $ns$ 才攒得出同样的相位"。

于是相位差与光程差的关系是 p15 的第一条式子：
\[
\delta=\frac{2\pi}{\lambda}\Delta L,\qquad \Delta L=L_2-L_1 .
\]
\aside{$\lambda$ 取\emph{真空}波长（本讲义记号表已约定）。若误用介质波长，$\delta$ 会差一个因子 $n$——这是干涉题里最常见的隐形错误。}

\subhdr{2.2 极大与极小：一条判据，两种写法}

p14 用相位写，p15 用光程写，是同一件事：

\tabtop
\noindent\begin{tabular}{@{}p{3.0cm}p{5.2cm}p{7.8cm}@{}}
\thead{条纹} & \thead{相位判据（p14）} & \thead{光程判据（p15）}\\
\tabmid
明纹（constructive） & $\delta=2m\pi$ & $\Delta L=m\lambda$\\
暗纹（destructive） & $\delta=(2m+1)\pi$ & $\Delta L=(2m+1)\dfrac{\lambda}{2}$\\
\tabbot

逐行读法：第一行把 $\delta=2m\pi$ 代进 $\delta=\frac{2\pi}{\lambda}\Delta L$，两边约掉 $2\pi/\lambda$ 就得 $\Delta L=m\lambda$；第二行同理。$m=0,\pm1,\pm2,\dots$，取负值对应屏上另一侧，不是"多出来的解"。

p15 的图（旋转双曲面族）说的是同一件事的几何版：等光程差的点集是一族以 $S_1,S_2$ 为焦点的旋转双曲面，它们与屏幕的交线才是明暗条纹。

\subhdr{2.3 一分钟小例}

一片 $n=1.5$、厚 $2\ \mu$m 的玻璃片，一束光垂直穿过它，光程比在空气中多走
\[
L-L_0=ns-s=(n-1)s=0.5\times2\ \mu\mathrm{m}=1\ \mu\mathrm{m}.
\]
这 $1\ \mu$m 就是 \S5 那道题的全部来源：盖不盖片，差别只在这一个 $(n-1)h$。

\flowerrule
% ---------------------------------------------------------------- §3
\hdr{\S3}{双缝几何：从 $R_2^2-R_1^2$ 到 $\Delta y=\lambda D/d$}

\subhdr{3.1 完整代数：p18 的 $R_2-R_1\approx\frac{d}{D}y$ 是怎么来的 \textbf{[补]}}

p17 已经把关键一步写出来了（"$R_2^2-R_1^2=2yd$"），p18 却直接跳到 $R_2-R_1\approx\frac{d}{D}y$。中间隔着一个容易被漏掉的因子 2，我们把它补全。

取缝 $S_1$ 在 $(0,\,d/2,\,D)$、$S_2$ 在 $(0,\,-d/2,\,D)$，屏上点 $P(x,y,0)$（p17 右侧的坐标就是这个约定）。则
\[
R_1^2=x^2+\Big(y-\frac{d}{2}\Big)^2+D^2,\qquad
R_2^2=x^2+\Big(y+\frac{d}{2}\Big)^2+D^2 .
\]
相减，$x^2$ 与 $D^2$ 抵消，只剩交叉项：
\[
R_2^2-R_1^2=2yd\quad\Longrightarrow\quad (R_2-R_1)(R_2+R_1)=2yd
\quad\Longrightarrow\quad R_2-R_1=\frac{2yd}{R_1+R_2}.
\]
现在做近似。当 $y\ll D$、$d\ll D$ 时，
\[
R_{1,2}=D\sqrt{1+\frac{(y\mp d/2)^2}{D^2}}\approx D\Big[1+\frac{(y\mp d/2)^2}{2D^2}\Big],
\]
两式相加：
\[
R_1+R_2\approx 2D\Big[1+\frac{y^2+d^2/4}{2D^2}\Big]=2D(1+\varepsilon),\qquad \varepsilon=\frac{y^2+d^2/4}{2D^2}.
\]
代回：
\[
\Delta L=R_2-R_1=\frac{2yd}{2D(1+\varepsilon)}\approx\frac{yd}{D}\quad(\varepsilon\ll1).
\]
\aside{$\varepsilon$ 就是这条近似的相对误差。量级：$D=1$ m、$y=5$ mm、$d=0.2$ mm 时 $\varepsilon\approx1.3\times10^{-5}$，完全可忽略。注意分子里的 $2yd$ \emph{不}约掉——它和分母里的 $2D$ 恰好抵消，这才留下 $yd/D$；只记住结果而不走这一步的人，常写出 $yd/(2D)$ 或 $2yd/D$。}

\subhdr{3.2 条纹位置与间距}

明纹条件 $\Delta L=m\lambda$ 与上式联立：
\[
\frac{y_m d}{D}=m\lambda\quad\Longrightarrow\quad y_m=m\frac{\lambda D}{d},
\]
相邻两条明纹（$m$ 与 $m+1$）的距离就是 p18 的方框式：
\[
\boxed{\;\Delta y=\frac{\lambda D}{d}\;}
\]
$m=0$ 给出 $y=0$，即屏幕中心 $O$ 点是零级明纹——这一点要单独说，因为它是 \S5 里"零级搬到哪里"的出发点。

\subhdr{3.3 图景}

\begin{center}
\begin{tikzpicture}[x=1cm,y=1cm,>={Stealth[length=1.8mm]},font=\small]
  % 源平面
  \draw[thick] (0,-2.3)--(0,2.3);
  \fill (0,0) circle (1.3pt) node[left=3pt]{$S$};
  % 缝平面（两条缝=两段空隙）
  \draw[thick] (6,-2.3)--(6,-0.52);
  \draw[thick] (6,-0.28)--(6,0.28);
  \draw[thick] (6,0.52)--(6,2.3);
  \fill (6,0.40) circle (1.2pt) node[above left=-1pt]{$S_1$};
  \fill (6,-0.40) circle (1.2pt) node[below left=-1pt]{$S_2$};
  % 光线
  \draw[->] (0,0)--(6,0.40);
  \draw[->] (0,0)--(6,-0.40);
  \draw[->] (6,0.40)--(12,1.45);
  \draw[->] (6,-0.40)--(12,1.45);
  \node at (7.4,0.95) {\scriptsize $R_1$};
  \node at (7.4,-0.72) {\scriptsize $R_2$};
  % 屏
  \draw[thick] (12,-2.3)--(12,2.3);
  \fill (12,1.45) circle (1.3pt) node[right=3pt]{$P$};
  \node at (12,0) [right=3pt] {$O$};
  \draw[->] (12.0,-0.15)--(12.0,-2.0) node[right=2pt] {\scriptsize $y$};
  % d 标注
  \draw[<->] (6.9,0.40)--(6.9,-0.40) node[midway,right=2pt] {\scriptsize $d$};
  % D 标注
  \draw[<->] (6,-1.95)--(12,-1.95) node[midway,above=1pt] {\scriptsize $D$};
\end{tikzpicture}
\end{center}

\subhdr{3.4 这套公式的适用边界（p21 的"空间周期性"要说的就是这个）}

p21 的原话是："The interference fringes are the macroscopic manifestation of the spatial periodicity of light waves." 翻译成人话：光的空间周期是 $\lambda$，装置把它放大成肉眼看得见的 $\Delta y$——放大倍数就是 $D/d$。

精确定义下，明纹条件是 $d\sin\theta=m\lambda$（$\theta$ 为 $P$ 相对中心的张角），而 $y=D\tan\theta$。改用 $\theta$ 求间距 \textbf{[补]}：
\[
\sin\theta=\frac{m\lambda}{d}\ \Longrightarrow\ \frac{\mathrm{d}\theta}{\mathrm{d}m}=\frac{\lambda}{d\cos\theta},
\qquad
\frac{\mathrm{d}y}{\mathrm{d}m}=D\sec^2\theta\cdot\frac{\mathrm{d}\theta}{\mathrm{d}m}
=\frac{\lambda D}{d\cos^3\theta}.
\]
所以严格地说，\emph{越往屏边缘，条纹越宽}（$\cos\theta<1$），等间距只是 $y\ll D$ 的下近似。这也解释了为什么"用双缝测波长"只在近轴小范围内成立。

\flowerrule
% ---------------------------------------------------------------- §4
\qbreak
\hdr{\S4}{例 5.1：用双缝测波长（课件 p19，全解）}

\qstem{（课件 p19，Example 5.1）杨氏双缝装置中，两缝间距为 $0.2$ mm，缝到屏的距离为 $1$ m。若第 1 条明纹与第 4 条明纹相距 $7.5$ mm，求光的波长。}

\subhdr{考点}
$\Delta y=\lambda D/d$ 的反用：把"两条纹的距离"换成"条纹间距"时，要数的是\emph{间隔数}，不是条纹数。

\subhdr{关键一步}
第 1 条到第 4 条之间隔了 $4-1=3$ 个间距（课件 p19 的解答里写的就是 $\Delta y=\frac{7.5}{4-1}=2.5$ mm）。

\subhdr{详解}
\[
\Delta y=\frac{7.5\ \mathrm{mm}}{4-1}=2.5\ \mathrm{mm},
\qquad
\lambda=\frac{d\,\Delta y}{D}=\frac{0.2\ \mathrm{mm}\times2.5\ \mathrm{mm}}{1\ \mathrm{m}} .
\]
统一到米：$0.2\ \mathrm{mm}=2\times10^{-4}$ m，$2.5\ \mathrm{mm}=2.5\times10^{-3}$ m，
\[
\lambda=\frac{2\times10^{-4}\times2.5\times10^{-3}}{1}=5.0\times10^{-7}\ \mathrm{m}=500\ \mathrm{nm}.
\]

\subhdr{答案}
$\lambda=500$ nm（与课件 p19 方框一致）。

\subhdr{自检}
量纲 $[\mathrm{m}][\mathrm{m}]/[\mathrm{m}]=\mathrm{m}$，对；$500$ nm 落在可见光内，对；代回 $y_4-y_1=3\lambda D/d=3\times500\times10^{-9}\times1/2\times10^{-4}=7.5$ mm，与题给一致。

\flowerrule
% ---------------------------------------------------------------- §5
\qbreak
\hdr{\S5}{例 5.2：云母片盖住一条缝（课件 p20，作业 HOME WORK TWO，全解）}

\qstem{（课件 p20，Ex.\,5.2 / HOME WORK TWO）杨氏双缝装置中，在上缝前放一片折射率 $n=1.58$ 的云母片。中央明纹（零级）移到了原来第 8 条明纹的位置。若光波长为 $550$ nm，求云母片的厚度。}

\subhdr{考点}
"插入一块介质" $=$ 给其中一束光加上一个\emph{恒定的}光程差。整道题只有一句话值得记：
\[
(n-1)h \quad\text{就是那片云母片贡献的额外光程。}
\]

\subhdr{关键一步}
在原来第 8 条明纹处 $P$，旧的光程差是 $8\lambda$；盖上片之后，这一处的光程差变成
\[
\Delta L_P'=8\lambda-(n-1)h .
\]
而 $P$ 现在是\emph{新的零级}，零级的定义就是光程差为零。于是
\[
8\lambda-(n-1)h=0\quad\Longrightarrow\quad h=\frac{8\lambda}{n-1}.
\]

\subhdr{详解：三条互相独立的路线}

\textbf{路线 A（光程差法，最直白）}：同上式，直接解出
\[
h=\frac{8\lambda}{n-1}=\frac{8\times550\ \mathrm{nm}}{0.58}=7586\ \mathrm{nm}\approx7.59\ \mu\mathrm{m}.
\]

\textbf{路线 B（级次搬运法）}：整幅条纹图样是"平移"的，不是"变形"的。屏上任一点的级次从 $m$ 变成
\[
m'=m-\frac{(n-1)h}{\lambda},
\]
因为光程差整体少了 $(n-1)h$，而每少一个 $\lambda$ 级次就减 1。零级现在落在旧的第 8 级处，即 $m=8$ 的点满足 $m'=0$：
\[
\frac{(n-1)h}{\lambda}=8\quad\Longrightarrow\quad h=\frac{8\lambda}{n-1}.
\]
"移动了 8 条"与"(n-1)h=8\lambda"是同一句话的两种说法：移动的条纹数 $\dfrac{(n-1)h}{\lambda}$ 可以不是整数——非整数表示新的零级恰好落在两条明纹之间。

\textbf{路线 C（相位法，你草稿上的写法）}：$P$ 处原来 $\delta_P=\frac{2\pi}{\lambda}\cdot8\lambda=16\pi$；盖上片后
\[
\delta_P'=16\pi-\frac{2\pi}{\lambda}(n-1)h=0\quad\Longrightarrow\quad (n-1)h=8\lambda .
\]

三条路线用的是同一个物理事实（额外光程 $(n-1)h$），只是记账方式不同：A 记在光程差上，B 记在级次上，C 记在相位上。考场上任选一条，写清楚"零级 $\Leftrightarrow$ 光程差为零"就不会错。

\subhdr{方向：条纹往哪边移？}

\emph{向被盖住的那条缝一侧移动。}本题盖的是上缝 $S_1$，所以条纹整体上移。

理由：盖片让上缝那束光"变慢"（光程变长），要重新凑成零级，就得让下缝那束光的几何路程补回来——观察点必须走到$\ $离上缝更近、离下缝更远的地方。用 \S3 的 $\Delta L=yd/D$ 写出来更干脆：新零级满足
\[
\frac{yd}{D}=(n-1)h>0\quad\Longrightarrow\quad y>0 ,
\]
即 $y$ 为正（上侧）。这也正好对上 p20 的图：$P$ 画在 $O$ 的上方。

\subhdr{答案}
\[
h=\frac{8\times550\ \mathrm{nm}}{1.58-1}=7.59\times10^{-6}\ \mathrm{m}\approx7.6\ \mu\mathrm{m}.
\]

有效数字的说明：$n=1.58$ 给到三位，$550$ nm 与"8 条"都只算两三位，答 $7.59\ \mu$m 或 $7.6\ \mu$m 都站得住；写 $7.6\ \mu$m 与"移动 8 条"这种整数级次相称。若非要一位不多地写，就写 $7.59\ \mu$m 并注明 $\lambda$ 视为三位。

\subhdr{与你手写稿的逐行对照}

你草稿的物理是对的，三处判断都对：$O$ 点原来 $\Delta\varphi_O=0$；$P$ 点原来 $\Delta\varphi_P=8k\lambda$（即 $16\pi$）；盖上片后 $\Delta\varphi_P'=\Delta\varphi_P-kd(n-1)$，令其为零。需要改的只有写法：

\begin{itemize}
\item \textbf{"$8k\lambda$"这个记法要拆开。}你用的 $k=2\pi/\lambda$（与课件 p13、p17 一致），所以 $8k\lambda=8\cdot2\pi=16\pi$ 数值没错；但它把"光程差 $8\lambda$"和"相位 $16\pi$"压进一个符号里，回头自己都容易看错。建议写成两行：先 $\Delta L_P=8\lambda$，再 $\delta_P=\frac{2\pi}{\lambda}\Delta L_P=16\pi$。
\item \textbf{别让 $k$ 兼职。}中文教材里 $k$ 常被用来表示\emph{级次}（"第 $k$ 级"），而你这里 $k$ 是波数。同一张纸上两个含义，考场上最容易翻车。级次请一律用 $m$。
\item \textbf{两个 $d$ 要分开。}你最后一行 $d=8\lambda/(n-1)$ 的 $d$ 是片厚，而 $\Delta L\approx yd/D$ 里的 $d$ 是缝距。交作业时写成 $h$，老师一眼就懂你在算什么。
\item \textbf{那个涂改处}正是"零级条件 $=0$"这一步。先在纸边写一句"新零级：$\Delta L=0$"，再往里代，就不用涂。
\end{itemize}

\subhdr{自检}
代回：$(n-1)h/\lambda=0.58\times7586\ \mathrm{nm}/550\ \mathrm{nm}=8.000$，对；量纲 $h=\lambda/(n-1)$ 是长度，对；数量级：8 个波长除以 $0.58$ 约 14 个波长，约 $7.6\ \mu$m，与"薄膜"常识相符。

\subhdr{同型变式（换数，供你确认"会迁移"）}
若改盖下缝、移动 5 条、$\lambda=600$ nm、$n=1.5$：$h=5\times600/0.5=6.0\ \mu$m，条纹整体下移。若反过来问"移动了几条"：条数 $=(n-1)h/\lambda$，允许非整数。

\flowerrule
% ---------------------------------------------------------------- §6
\hdr{\S6}{条纹对比度：为什么有时条纹"看不清"}

\subhdr{6.1 定义与来源（p22--p24）}

p22 给出对比度（可见度）
\[
\gamma=\frac{I_{\max}-I_{\min}}{I_{\max}+I_{\min}} .
\]
把 $I=I_1+I_2+2\sqrt{I_1I_2}\cos\delta$ 的 $\cos\delta=\pm1$ 代进去：
\[
I_{\max}=I_1+I_2+2\sqrt{I_1I_2},\qquad I_{\min}=I_1+I_2-2\sqrt{I_1I_2}.
\]
于是
\[
\gamma=\frac{2\sqrt{I_1I_2}}{I_1+I_2}=\frac{2A_1A_2}{A_1^2+A_2^2}=\frac{2\rho}{1+\rho^2},\qquad \rho\equiv\frac{A_1}{A_2}.
\]
\aside{这里用 $I=A^2/2$ 消掉了时间平均的 $1/2$，所以 $\gamma$ 只依赖\emph{振幅比} $\rho$，与总光强无关。这就是为什么"把灯调亮"不会改善对比度——它只让条纹更亮，明暗差别按比例一起变亮。}

p23、p24 把结果整理成
\[
I=I_0(1+\gamma\cos\delta),\qquad I_0=I_1+I_2 .
\]
p24 页脚小字说明了 $I_0$ 的含义：两束光"各管各"时的总强度（即把交叉项拿掉之后的和）。

\subhdr{6.2 两个极限与"非相干"的数学含义}

$\rho=1$（两束等强）：$\gamma=1$，$I_{\max}=4I_0/2$ 型的完全调制；$\rho\to0$ 或 $\rho\to\infty$：$\gamma\to0$，条纹消失。

p23 的原话是 "When $\gamma=0$, it degenerates to the incoherent (non-interfering) sum"。数学上这就是 \S1.2 里"交叉项被平均掉"的那一档：$\gamma=0$ 时 $I=I_0=I_1+I_2$，强度在屏上是常数，没有条纹——不是"条纹很弱"，而是"根本没有条纹"。

\subhdr{6.3 常数报户口}

$\gamma$ 里出现的三个数各有出处：交叉项前的 $\boxed{2}$ 来自 \S1.2 展开式的 $2E_1E_2$；分子 $2A_1A_2$ 来自 $I_{\max}-I_{\min}=4\sqrt{I_1I_2}$ 与分母 $2(I_1+I_2)$ 相除；$I=A^2/2$ 里的 $1/2$ 是时间平均（\S1.2），在 $\gamma$ 里被约掉。三者分工不同，不要混着背。

\subhdr{6.4 一分钟小例}
若 $A_2=3A_1$，则 $\rho=1/3$，
\[
\gamma=\frac{2/3}{1+1/9}=0.6 .
\]
验算：$I_{\max}:I_{\min}=(A_1+3A_1)^2:(3A_1-A_1)^2=16:4=4:1$，故 $\gamma=\frac{4-1}{4+1}=0.6$，一致。

\flowerrule
% ---------------------------------------------------------------- §7
\hdr{\S7}{波前分割的代价：p25 那个"代价呢？？？" \textbf{[补]}}

p25 的绿框中文原话是：\textbf{"杨氏双缝干涉的独特优势就是获得两个相位稳定的子波源 $s_1$ 和 $s_2$，且无视初始相位项，代价呢？？？"}——课件在这一页提了问，下一页直接跳到双棱镜，没有回答。这一节替它回答。

\subhdr{7.1 机制：扩展光源的每一"点"各画一套条纹}

p25 的 ⚠ 框讲的是现象：不能用普通灯或扩展光源，必须用点光源或极窄的缝。物理原因是——
\begin{itemize}
\item 扩展光源上不同的点，到两缝的距离差不同（多出 $\xi d/l$，$\xi$ 是源点离轴高度，$l$ 是源到缝的距离）；
\item 每个点各自在屏上生出一套条纹，彼此\emph{错开}；
\item 许多套错开的条纹非相干叠加，对比度被"抹平"（p25 原话 "the fringe contrast to drop sharply and eventually wash out"）。
\end{itemize}

\subhdr{7.2 定量：平移了几条？}

源点高 $\xi$ 时，多出的光程差 $\Delta L_{\text{源}}=\xi d/l$；它让整套条纹沿屏移动。用 $\Delta L=yd/D$ 反解移动量：
\[
\frac{yd}{D}=\frac{\xi d}{l}\quad\Longrightarrow\quad y=\frac{\xi D}{l},
\]
换成"条纹数"：
\[
\frac{y}{\Delta y}=\frac{\xi D/l}{\lambda D/d}=\frac{\xi d}{\lambda l}.
\]

\subhdr{7.3 宽度 $b$ 的缝光源：对比度按 sinc 衰减}

沿 $\xi$ 均匀取平均（宽度 $b$，令 $u\equiv\dfrac{\pi b d}{\lambda l}$）：
\[
\frac1b\int_{-b/2}^{b/2}\cos\Big(\delta+\frac{2\pi\xi d}{\lambda l}\Big)\mathrm{d}\xi
=\cos\delta\cdot\frac{\sin u}{u}.
\]
于是
\[
\gamma=\Big|\frac{\sin u}{u}\Big|,\qquad u=\frac{\pi b d}{\lambda l}.
\]
第一个零点在 $u=\pi$，即
\[
b_{\text{c}}=\frac{\lambda l}{d}.
\]
量级例子：$l=0.5$ m、$d=0.2$ mm、$\lambda=550$ nm，得 $b_{\text{c}}=1.375$ mm——单缝宽度超过约 $1.4$ mm，条纹就基本看不见了。

\subhdr{7.4 这就是"代价"}

为了"两个子波源相位稳定"，必须\emph{牺牲光通量}：把广阔的光源挡到只剩一条极窄的缝，只用两个小区域参与干涉。相干性越好，能用的光越少——这是分波前干涉（以及后来所有干涉仪）都要付的代价。

\flowerrule
% ---------------------------------------------------------------- §8
\hdr{\S8}{三种分波前装置：同一个模板，三种"造双源"的办法}

\subhdr{8.1 统一模板}

p26--p28 三个装置看着不同，做法完全一样：\emph{用镜子/棱镜把同一束光"折"成两个虚像，当成一双相干点源，然后回到 \S3}。只要求出等效双源间距 $d_{\text{eq}}$，条纹间距就是
\[
\Delta x=\frac{\lambda\,(B+C)}{d_{\text{eq}}} ,
\]
其中 $B+C$ 是"虚源面到屏"的距离（双镜、双棱镜用 $B+C$；洛埃镜用 $D$）。

\subhdr{8.2 三个装置对照}

\tabtop
\noindent\begin{tabular}{@{}p{2.4cm}p{3.4cm}p{5.0cm}p{5.0cm}@{}}
\thead{装置（页）} & \thead{等效双源间距} & \thead{条纹间距} & \thead{几何来源}\\
\tabmid
双平面镜（p26） & $d\approx2\alpha B$ & $\Delta x=\dfrac{(B+C)\lambda}{2\alpha B}$ & 两虚像关于两镜面对称，对交点 $O$ 的张角为 $2\alpha$\\
双棱镜（p27） & $d\approx2(n-1)\alpha B$ & $\Delta x=\dfrac{(B+C)\lambda}{2(n-1)\alpha B}$ & 上下两半各偏折 $(n-1)\alpha$\\
洛埃镜（p28） & $d=2a$ & $\Delta x=\dfrac{D\lambda}{2a}$ & 源与其镜像关于镜面对称，相距 $2a$\\
\tabbot

逐行读法：
\begin{itemize}
\item 第一行：$S$ 到镜交点的距离为 $B$，两镜夹角 $\alpha$；镜面反射把 $S$ 分别映成 $S_1$、$S_2$，两虚像各距交点 $B$，夹角为 $2\alpha$，故弦长 $d\approx2B\sin\alpha\approx2\alpha B$（小角）。
\item 第二行：双棱镜顶角 $\alpha$ 很小，每一半像一个顶角 $\alpha$ 的薄棱镜，把光偏折 $(n-1)\alpha$（推导见旁注），两虚像夹角 $2(n-1)\alpha$，故 $d\approx2(n-1)\alpha B$。
\item 第三行：只有一个平面镜，$S$ 的镜像与 $S$ 关于镜面对称，间距是高度的两倍 $2a$；屏到源面的距离记作 $D$，于是 $\Delta x=D\lambda/(2a)$。
\end{itemize}

\aside{薄棱镜偏向角 $\delta_d=(n-1)\alpha$ 的一行推导：小角下 $\sin\theta\approx\theta$，两次折射给 $i_1=nr_1$、$i_2=nr_2$，而棱镜内两法线夹角满足 $r_1+r_2=\alpha$，故偏向角 $\delta_d=(i_1+i_2)-\alpha=n\alpha-\alpha=(n-1)\alpha$。}

\subhdr{8.3 洛埃镜的半波损失（p28）}

p28 右侧粉框原话："At grazing incidence, the reflected light undergoes a phase reversal of $\pi$ (half-wave loss)." 为什么？

反射发生在空气 $\to$ 玻璃的界面上（光疏 $\to$ 光密）。正入射时振幅反射系数
\[
r=\frac{n_1-n_2}{n_1+n_2}<0\quad(n_2>n_1),
\]
负号就是 $\pi$ 的相位突变。掠入射时，垂直于入射面的分量在任何角度都反相；平行分量在掠入射时也趋近反相——所以 p28 说"掠入射时"发生 $\pi$ 相移，是对的。

后果：贴着镜面那一点（几何光程差为 0）出现的是\textbf{暗纹}，而不是杨氏那样的明纹。这就是 p28 粉框里"bright and dark fringes are reversed"的意思，也是洛埃镜被用来"验证半波损失"的原因。

\subhdr{8.4 一分钟小例（双棱镜）}
$\alpha=0.5^\circ$、$n=1.5$、$B=0.2$ m、$C=1.0$ m、$\lambda=550$ nm：
\[
d\approx2\times0.5\times0.008727\times0.2=1.745\ \mathrm{mm},\qquad
\Delta x=\frac{1.2\times550\times10^{-9}}{1.745\times10^{-3}}=0.378\ \mathrm{mm}.
\]

\subhdr{8.5 课件这两页的一个错位要说明}
p26 的标题写的是 "Fresnel Biprism"，但正文第一条是 "Made of two plane mirrors in contact, forming a small angle $\alpha$"——图与文都是\textbf{双平面镜}；真正的\textbf{双棱镜}在 p27。两者公式只差一个 $(n-1)$，但物理装置不同（一个靠反射，一个靠折射）。看课件时别把 p26 的 $\alpha$ 当成棱镜顶角去套 p27 的式子。

\flowerrule
% ---------------------------------------------------------------- §9
\hdr{\S9}{公式卡与四步法}

\subhdr{9.1 公式卡}

\tabtop
\noindent\begin{tabular}{@{}p{5.4cm}p{5.6cm}p{5.0cm}@{}}
\thead{公式} & \thead{成立条件} & \thead{用在哪}\\
\tabmid
$\delta=\dfrac{2\pi}{\lambda}\Delta L$ & 单色光；$\lambda$ 为真空波长 & 一切干涉题的换算枢纽\\
$\Delta L=m\lambda$；$\Delta L=(2m+1)\dfrac{\lambda}{2}$ & 两束相干光、等强或不等强均可 & 明纹 / 暗纹\\
$\Delta L\approx\dfrac{yd}{D}$ & $y\ll D$、$d\ll D$ & 双缝类装置的几何核\\
$\Delta y=\dfrac{\lambda D}{d}$ & 同上 & 条纹间距、测波长\\
$\Delta L\to\Delta L-(n-1)h$ & 薄片垂直（近垂直）穿过一束光 & 插片移条纹\\
$\gamma=\dfrac{2A_1A_2}{A_1^2+A_2^2}$，$I=I_0(1+\gamma\cos\delta)$ & 两束偏振方向平行、相位差稳定 & 条纹清晰度\\
$\Delta x=\dfrac{(B+C)\lambda}{2(n-1)\alpha B}$ & 小顶角薄棱镜 & 双棱镜\\
$\Delta x=\dfrac{D\lambda}{2a}$ & 掠入射、半波损失 & 洛埃镜\\
\tabbot

\subhdr{9.2 四步法}

\begin{enumerate}
\item \textbf{画图定位}：源——缝（或镜、棱镜）——屏，标出 $d$、$D$，并标清\emph{哪一束}路过了介质、走了多厚。
\item \textbf{写光程差}：几何段照写，介质段乘 $n$；把"插入薄片"翻译成"某束加 $(n-1)h$"。
\item \textbf{套判据}：先用 $\delta=\frac{2\pi}{\lambda}\Delta L$ 或直接明暗条件；要求零级就令 $\Delta L=0$；要求"移动几条"就数 $\Delta L$ 变了几个 $\lambda$。
\item \textbf{代数字 + 三查}：查量纲、查数量级（条纹间距通常是 mm 级、片厚是 $\mu$m 级）、查极限（两束等强时 $\gamma$ 是否为 1；$m=0$ 时是否回到中心）。
\end{enumerate}

\flowerrule
% ---------------------------------------------------------------- §10
\hdr{\S10}{易错点（每条都能指回某一页的某一步）}

\begin{enumerate}
\item \textbf{两个 $d$ 混用}（p16--p18 的缝距 vs p20 的片厚）。后果：公式对、字母错，自己复查时看不懂。对策：片厚一律写 $h$。
\item \textbf{几何近似漏因子}：由 $2yd/(R_1+R_2)$ 往 $yd/D$ 走时，分母的 $R_1+R_2\approx2D$ 与分子的 2 抵消；写成 $yd/(2D)$ 会让 $\Delta y$ 差一倍（\S3.1）。
\item \textbf{相位差少 $2\pi$}：写成 $\delta=\Delta L/\lambda$。用 $\delta=\frac{2\pi}{\lambda}\Delta L$（\S2.1）。
\item \textbf{移动方向反了}：向\emph{被盖住}的那条缝一侧（\S5）。
\item \textbf{"第 8 条"数成 $m=7$}：零级是第 1 条，故"第 8 条"对应 $m=8$。若按"把中央算第 1 条"的另一种计法，答案会变成 $6.6\ \mu$m；本题与 p20 的叙述一致取 $m=8$（即 $7.6\ \mu$m）。\textbf{交作业时写清你的计数约定}，老师才不会判错。
\item \textbf{以为"随便两只灯也能干涉"}（p8--p9）：初相位随机，$\delta(t)$ 漂移，交叉项平均为零。
\item \textbf{用扩展光源直接照双缝}（p25）：对比度按 sinc 塌掉（\S7）。
\item \textbf{洛埃镜忘了半波损失}（p28）：明暗整体反了，中心应是暗纹。
\item \textbf{对比度里拿强度比代振幅比}：$\gamma$ 的自变量是 $\rho=A_1/A_2$，不是 $I_1/I_2$（\S6.1）。
\end{enumerate}

\flowerrule
% ---------------------------------------------------------------- §11
\hdr{\S11}{缺口清单：课件这 16 页的跳步与补法}

\gapitem{代数跳步（p18）}{原文只写 "$R_2-R_1\approx\frac{d}{D}y$"，没有交代分母为什么是 $2D$、分子为什么是 $2yd$。本讲义 \S3.1 补出 $R_2^2-R_1^2=2yd\to R_2-R_1=2yd/(R_1+R_2)\to R_1+R_2\approx2D$ 的全链，并给出相对误差 $\varepsilon=(y^2+d^2/4)/(2D^2)$。}

\gapitem{机制缺口（p6）}{原文直接给 $I=I_1+I_2+2\sqrt{I_1I_2}\cos\delta$，跳过"对时间取平均"。本讲义 \S1.2 补 $\langle E^2\rangle=A^2/2$ 与两束光的交叉项展开，顺带说明频率不同时它为什么归零。}

\gapitem{换算缺口（p14 与 p15）}{p14 的极大极小用相位写（$2m\pi$、$(2m+1)\pi$），p15 用光程写（$m\lambda$、$(2m+1)\lambda/2$），两页之间没有写出 $2\pi/\lambda$ 这一步换算。本讲义 \S2.2 给出对照表与逐行读法。}

\gapitem{答案缺口（p20）}{作业页只有题面（"The central bright fringe (zero order) shifts to the position of the original 8th bright fringe"），没有给判据，也没有说条纹往哪边移。本讲义 \S5 给三条独立路线、方向论证、有效数字讨论，并对你手写稿的记号逐条核对。}

\gapitem{提问缺口（p25）}{原文用中文提问"代价呢？？？"，下一页直接进入双棱镜，没有回答。本讲义 \S7 补空间相干性的定量结果 $\gamma=|\sin u/u|$、$u=\pi bd/(\lambda l)$，并给出相干宽度 $b_c=\lambda l/d$。}

\gapitem{图文错位（p26）}{标题 "Fresnel Biprism" 与正文 "Made of two plane mirrors in contact" 不一致；图是双平面镜。本讲义 \S8 把双镜与双棱镜分开讲，并点明两者公式差 $(n-1)$。}

\gapitem{机制缺口（p28）}{原文只说 "at grazing incidence, the reflected light undergoes a phase reversal of $\pi$ (half-wave loss)"，没有说为什么掠入射也成立。本讲义 \S8.3 补反射系数 $r=(n_1-n_2)/(n_1+n_2)<0$ 的来源，并说明掠入射时两个分量的相位行为。}

\flowerrule
% ---------------------------------------------------------------- §12
\hdr{\S12}{三句话总结与自检}

\subhdr{12.1 三句话}

\begin{enumerate}
\item 干涉的全部内容是\emph{光程差}：$\Delta L$ 变一个 $\lambda$，条纹就挪一格；插入介质只是把某束光的 $\Delta L$ 加上 $(n-1)h$。
\item 双缝把 $\Delta L$ 变成可测的几何量 $\Delta L\approx yd/D$，于是 $\Delta y=\lambda D/d$——装置把光的空间周期 $\lambda$ 放大 $D/d$ 倍。
\item 能把条纹看清的前提是相位稳定与振幅相近：前者要求点光源（窄缝），后者要求两束强度可比（对比度 $\gamma$）。
\end{enumerate}

\subhdr{12.2 自检三题（答案在题后，自检用，不要当独立检验）}

\textbf{1.} 双缝 $d=0.5$ mm、$D=2$ m、$\lambda=600$ nm。求条纹间距与第 3 级明纹到中心的距离。

\textbf{2.} 一条缝前盖 $n=1.5$、$h=3\ \mu$m 的薄片，$\lambda=500$ nm。条纹移动几条？往哪边？

\textbf{3.} 双棱镜 $\alpha=1^\circ$、$n=1.5$、$B=0.25$ m、$C=1.5$ m、$\lambda=500$ nm。求条纹间距。

\medskip
\textbf{答案与要点：}
\begin{enumerate}
\item $\Delta y=\lambda D/d=600\times10^{-9}\times2/0.5\times10^{-3}=2.4$ mm；$y_3=3\Delta y=7.2$ mm。
\item 条数 $=(n-1)h/\lambda=0.5\times3\ \mu\mathrm{m}/500\ \mathrm{nm}=3$；向被盖住的那条缝一侧移动。零级恰好落在旧的第 3 级处。
\item $d\approx2(n-1)\alpha B=2\times0.5\times0.017453\times0.25=4.363$ mm；$\Delta x=(B+C)\lambda/d=1.75\times500\times10^{-9}/4.363\times10^{-3}\approx0.20$ mm。
\end{enumerate}

\flowerrule

\callout{\textbf{下一步。}把 \S5 的三条路线各写一遍（只用一句话概括每条的"记账对象"），然后闭卷把 p20 这道题完整写一遍：画图、写 $\Delta L$、令零级为零、代数字、自检。写的时候把片厚记作 $h$，把 $k$ 只留给波数。写完可以拍照发我，我只核对"零级条件"和"方向"这两处，其余不重复讲。}

\end{document}
